% Requirements by Unizar's press: 
%   1. Page size 18x25cm (B5)
%   2. All chapters begin on an odd page (meaning right-hand side of a book)
%   3. Figure caption must have a smaller font than the regular text

% Guidelines of the FMC department in addition to Unizar's press requirements:
%   1. Upper margin 2cm, lower margin 2cm, inner margin 3cm, outer margin 2cm -> textwidth = 13cm, textheight = 210 cm
%   2. Total width dedicated to header 1cm, which takes up 1cm of the 2cm dedicated to the upper margin

% GOALS: Comply with the requirements and guidelines. Format the text only once. Be able to print both in A4 and in 18x25cm
% SOLUTION: Always format text in a 13x(21+1)cm window but adjust margins to change the paper size depending on the purpose
% NOTE: Requirement 2 is automatically satisfied by using \documentclass{book}


% Margins, widths, heights and spacings
\usepackage{geometry} 
\geometry{includehead} % consider header as part of the text body when setting margins, width and height

\ifforpress
    \geometry{paperheight=25cm, paperwidth=18cm} % custom paper size of Unizar's press
    \geometry{inner=3cm, top=1cm, width=130mm, height=220mm} % margins width and height are set according to requirement 1 and guidelines 1 and 2
\else
    \geometry{a4paper} % typical A4 size
    \geometry{left=40mm, top=45mm, width=130mm, height=220mm} % margins are set so that right and left are symmetric, width and height according to requirement 1 and guidelines 1 and 2
\fi

% adjust total header width to 1cm = 10mm to comply with guidelines 1 and 2
\setlength{\headsep}{5mm} % spacing between header and tex
\setlength{\headheight}{5mm} % vertical width of header

\setlength{\parskip}{0.2cm} % spacing between paragraphs
\renewcommand{\baselinestretch}{1.05} % multiplier to \baselineskip which sets the distance between the bottom of two consecutive lines


% Fancy header style options
\usepackage{fancyhdr}

% the style of regular pages
\fancypagestyle{regular}{
    \fancyhead[LE, RO]{\small \rm \hyperlink{toc}{\thepage}} % page number, in the left on odd pages and in the right on even pages, with a hyperlink to the toc
    \fancyhead[RE]{\small \sl \leftmark} % chapter in the right on even pages
    \fancyhead[LO]{\small \sl \rightmark} % section in the left on odd pages
    \renewcommand{\chaptermark}[1]{\markboth{\chaptername \ \thechapter.\ ##1}{}} % I don't really understand this
    \renewcommand{\sectionmark}[1]{\markright{\thesection.\ ##1}{}} % I don't really understand this
    \renewcommand{\headrulewidth}{0.3pt} % width of the line underscoring the headers
}

% a special page style for unnumbered chapters (abstract, introduction, conclusions)
\fancypagestyle{unnumbered}{
    \fancyhead[LE, RO]{\small \rm \hyperlink{toc}{\thepage}} 
    \fancyhead[RE]{\small \sl \leftmark} 
    \fancyhead[LO]{\small \sl \rightmark} 
    \renewcommand{\headrulewidth}{0.3pt} 
}

% a special page style for the appendix
\fancypagestyle{appendix}{
    \fancyhead[LE, RO]{\small \rm \hyperlink{toc}{\thepage}} % page number, in the left on odd pages and in the right on even pages, with a hyperlink to the toc
    \fancyhead[RE]{\small \sl \leftmark} % chapter in the right on even pages
    \fancyhead[LO]{\small \sl \rightmark} % section in the left on odd pages
    \renewcommand{\chaptermark}[1]{\markboth{\appendixname \ \thechapter.\ ##1}{}} % style of how the chapter mark is displayed on headers, mainly to make it lowercase
    \renewcommand{\sectionmark}[1]{\markright{\thesection.\ ##1}{}} % style of how the section mark is displayed on headers, mainly to make it lowercase
    \renewcommand{\headrulewidth}{0.3pt} % width of the line underscoring the headers
}

% a special page style for the table of contents
\fancypagestyle{toc}{
    \fancyhf{} % clear all header and footer fields before defining custom ones, somehow neccesary to avoid page number at the bottom of the page
    \fancyhead[LE, RO]{\small \rm \thepage} % page number, in the left on even pages and in the right on odd pages
    \fancyhead[RE, LO]{\small \sl Contents}   % "Contents" in the right on even pagesand in the left on odd pages
    \renewcommand{\headrulewidth}{0.3pt} % width of the line underscoring the headers
}

% a special page style for the bibliography, for some reason the unnumbered style does not work properly
\fancypagestyle{biblio}{
    \fancyhead[LE, RO]{\small \rm \hyperlink{toc}{\thepage}} % page number, in the left on even pages and in the right on odd pages, with a hyperlink to the toc
    \fancyhead[RE, LO]{\small \sl Bibliography} % "Bibliography" in the right on even pages and in the left on odd pages
    \renewcommand{\headrulewidth}{0.3pt} % width of the line underscoring the headers
}

% redefinition of the plain style that is invoked automatically for the first page of parts and chapters
\fancypagestyle{plain}{
    \fancyhf{} % clear all header and footer fields before defining custom ones, somehow neccesary to avoid page number at the bottom of the page
    \renewcommand{\headrulewidth}{0.pt} % eliminate the header line
}


% Remove headers (and everything else) from supposedly emtpy pages, like empty verso pages between chapters
\usepackage{emptypage}

% To import pages from a PDF file (for the standard cover of the submitted version)
\usepackage{pdfpages}

% For displaying source code
\usepackage{listings}

% To improve the subconscius appeal of the text through typographic microadjustments and reduce the number of underfull and overfull badness warnings
\usepackage{microtype}

% To have a more modern font encoding
\usepackage[T1]{fontenc}

% To adjust spacings in the table of contents
\usepackage[titles]{tocloft}
\setlength{\cftparskip}{0.1cm} % to adjust vertical spacing between entries
\cftsetrmarg{3.55em plus 1fil} % to adjust marging between entry and page number. default: 2.55em

% Figures and math
\usepackage{graphicx}
\usepackage{amstext, amsmath, amssymb, physics, mathtools}
\usepackage[small]{caption} % make caption font size small according to requirement 3

% Hyperrefs
\ifforpress
    \usepackage[colorlinks=false, hidelinks=true]{hyperref} % turn off colored links to reduce printing cost (make most pages use black ink only)
\else
    \usepackage[colorlinks=true, allcolors=blue]{hyperref} % links in blue otherwise
\fi

% Bibliography
\usepackage[style=phys, % aps style
citestyle=numeric-comp, % shorten multiple citations as 1-3 instead of 1, 2, 3
biblabel=brackets, % bib entries are indexed with [n] instead of superscript n in biliography
sorting=none, % bib entries appear in order of citation
backref=true, % add backreference in bib entry to citation
backrefstyle=three % style of backreferences
]{biblatex} % biblatex uses the biber engine (instead of the bibtex engine) by default

\usepackage{bibentry} % To include bibliographic entries in regular text (for the list of publications)
